\documentclass[12pt]{article}
\usepackage[ukrainian,shorthands=off]{babel}   % shorthands=off: символ " не активний (потрібно для міток "..." у TikZ всередині \zadtask)
\usepackage{fontspec}
% --- НАЛАШТУВАННЯ ШРИФТІВ ---
% Шрифти Schoolbook-*.otf лежать у корені проєкту. Шлях підбирається автоматично:
% ../ (компіляція з папки теми), ./ (корінь проєкту), ../../ (підпапки теми 28); інакше -- TeX Gyre Schola.
\newcommand{\nmtSchoolbook}[1]{\setmainfont{Schoolbook-Regular.otf}[
    Path = #1,
    BoldFont = Schoolbook-Bold.otf,
    ItalicFont = Schoolbook-Italic.otf,
    BoldItalicFont = Schoolbook-BoldItalic.otf
]}
\IfFileExists{../Schoolbook-Regular.otf}{\nmtSchoolbook{../}}{%
 \IfFileExists{./Schoolbook-Regular.otf}{\nmtSchoolbook{./}}{%
  \IfFileExists{../../Schoolbook-Regular.otf}{\nmtSchoolbook{../../}}{%
   \setmainfont{texgyreschola}[Extension=.otf, UprightFont=*-regular, BoldFont=*-bold, ItalicFont=*-italic, BoldItalicFont=*-bolditalic]}}}
\usepackage[italic]{mathastext}
\usepackage[a4paper,margin=1.8cm,bottom=2cm,top=2cm]{geometry}
\usepackage{amsmath,amssymb}
\usepackage{tikz}
\usepackage{pgfplots}
\pgfplotsset{compat=1.18}
\usepgfplotslibrary{fillbetween}
\usetikzlibrary{calc,patterns,angles,quotes,intersections,babel,3d,shapes.symbols,shapes.geometric,decorations.pathreplacing,shadings,arrows.meta,hobby}
\usepackage{xcolor,array,fancyhdr,enumitem,multicol}
\usepackage{colortbl,multirow,diagbox,needspace}
\usepackage{tcolorbox}
\tcbuselibrary{skins,breakable}
\usepackage[hidelinks]{hyperref}
% водяний знак; щоб зібрати без нього (версія для вчителів), перед \documentclass пишуть \def\nmtnowatermark{}
\ifdefined\nmtnowatermark\else
\usepackage[firstpage=false, text=@pvtr2525, scale=0.6, color=gray!12]{draftwatermark}
\fi

% --- АВТОСТИСКАННЯ: рисунок/сітка, ширша за колонку, масштабується до ширини колонки ---
\newsavebox{\nmtfitbox}
\newenvironment{nmtfit}{\begin{lrbox}{\nmtfitbox}}{\end{lrbox}%
  \ifdim\wd\nmtfitbox>\linewidth \resizebox{\linewidth}{!}{\usebox{\nmtfitbox}}\else\usebox{\nmtfitbox}\fi}
\let\nmtIncludegraphicsOrig\includegraphics
\renewcommand{\includegraphics}[2][]{\begin{nmtfit}\nmtIncludegraphicsOrig[#1]{#2}\end{nmtfit}}


\definecolor{mainGreen}{RGB}{34, 120, 64}
\definecolor{yearOrange}{RGB}{220, 100, 30}
\definecolor{yearcolor}{RGB}{220, 100, 30}
\definecolor{titleBg}{RGB}{225, 240, 220}
\definecolor{instrBg}{RGB}{255, 240, 220}
\definecolor{instrBorder}{RGB}{220, 150, 80}
\definecolor{headerblue}{RGB}{0, 102, 204}   % використовується в рисунках старої бази

\setlength{\parindent}{0pt}
\emergencystretch=1.5em

\pagestyle{fancy}
\fancyhf{}
\renewcommand{\headrulewidth}{0.4pt}
\renewcommand{\footrulewidth}{0.4pt}
\fancyhead[C]{\small\color{gray!80} Варіант № 4}
\fancyhead[R]{\small\color{mainGreen}\textbf{\href{https://t.me/pvtr2525}{@pvtr2525}}}
\fancyfoot[L]{\small\color{gray!80}\href{https://t.me/pvtr2525}{ТГ: @pvtr2525}}
\fancyfoot[C]{\small\color{gray!80}\thepage}
\fancyfoot[R]{\small\color{gray!80}НМТ з математики}

% --- ТАБЛИЦІ ВІДПОВІДЕЙ ---
% ширина клітинки = (ширина рядка - відступи) / 5: на всю сторінку це ~3 см, у minipage поруч із рисунком таблиця стискається
\newlength{\nmtcellw}
\newcommand{\nmtSetCell}{\setlength{\nmtcellw}{\dimexpr(\linewidth-12\tabcolsep-6\arrayrulewidth)/5\relax}}
\newcommand{\answerTable}[5]{%
\par\nopagebreak\vspace{0.15cm}
\noindent\nmtSetCell
\begin{tabular}{|*{5}{>{\centering\arraybackslash}m{\nmtcellw}|}}
\hline
\rule[-0.15cm]{0pt}{0.6cm}\textbf{А} & \rule[-0.15cm]{0pt}{0.6cm}\textbf{Б} & \rule[-0.15cm]{0pt}{0.6cm}\textbf{В} & \rule[-0.15cm]{0pt}{0.6cm}\textbf{Г} & \rule[-0.15cm]{0pt}{0.6cm}\textbf{Д} \\
\hline
\rule[-0.2cm]{0pt}{0.75cm}#1 & \rule[-0.2cm]{0pt}{0.75cm}#2 & \rule[-0.2cm]{0pt}{0.75cm}#3 & \rule[-0.2cm]{0pt}{0.75cm}#4 & \rule[-0.2cm]{0pt}{0.75cm}#5 \\
\hline
\end{tabular}
\par\vspace{0.25cm}
}
\newcommand{\answerTableTall}[5]{%
\par\nopagebreak\vspace{0.15cm}
\noindent\nmtSetCell
\begin{tabular}{|*{5}{>{\centering\arraybackslash}m{\nmtcellw}|}}
\hline
\rule[-0.2cm]{0pt}{0.7cm}\textbf{А} & \rule[-0.2cm]{0pt}{0.7cm}\textbf{Б} & \rule[-0.2cm]{0pt}{0.7cm}\textbf{В} & \rule[-0.2cm]{0pt}{0.7cm}\textbf{Г} & \rule[-0.2cm]{0pt}{0.7cm}\textbf{Д} \\
\hline
\rule[-0.45cm]{0pt}{1.2cm}#1 & \rule[-0.45cm]{0pt}{1.2cm}#2 & \rule[-0.45cm]{0pt}{1.2cm}#3 & \rule[-0.45cm]{0pt}{1.2cm}#4 & \rule[-0.45cm]{0pt}{1.2cm}#5 \\
\hline
\end{tabular}
\par\vspace{0.25cm}
}
\newcommand{\instructionBox}[1]{%
\par\vspace{0.3cm}
\begin{tcolorbox}[colback=instrBg, colframe=instrBorder, boxrule=0.6pt, arc=2pt,
    left=8pt, right=8pt, top=4pt, bottom=4pt]
\centering\small\bfseries #1
\end{tcolorbox}
\par\vspace{0.15cm}
}
\newcommand{\sectionTitle}[1]{%
\par\vspace{0.4cm}
\begin{tcolorbox}[colback=titleBg, colframe=titleBg, boxrule=0pt, arc=8pt,
    halign=center, fontupper=\Large\bfseries\color{mainGreen}]
\hyphenpenalty=10000 \exhyphenpenalty=10000 #1
\end{tcolorbox}
\par\vspace{0.3cm}
}
\newcommand{\task}[2]{\noindent\textbf{#1.}\ \ #2\par\vspace{0.2cm}}
% --- НМТ-СТИЛЬ ПОЛЕ ВІДПОВІДІ ---
\newcommand{\nmtAnswerBox}{%
\par\nopagebreak\vspace{0.2cm}\noindent
Відповідь:\ \
\begin{tabular}{@{}*{4}{|>{\centering\arraybackslash}p{0.8cm}}|@{\hspace{0.1cm}\textbf{,}\hspace{0.1cm}}*{3}{|>{\centering\arraybackslash}p{0.8cm}}|@{}}
\hline
\rule[-0.2cm]{0pt}{0.7cm} & & & & & & \\
\hline
\end{tabular}
\par\vspace{0.3cm}
}
% --- СІТКА ДЛЯ МАТЧИНГУ ---
\newsavebox{\nmtgridbox}
\newcommand{\matchingGrid}{%
    \sbox{\nmtgridbox}{\matchingGridRaw}%
    \ifdim\wd\nmtgridbox>\linewidth \resizebox{\linewidth}{!}{\usebox{\nmtgridbox}}\else\usebox{\nmtgridbox}\fi}
\newcommand{\matchingGridRaw}{%
    \begingroup
    \renewcommand{\arraystretch}{1.15}
    \setlength{\tabcolsep}{4pt}
    \footnotesize
    \begin{tabular}{r|c|c|c|c|c|}
         \multicolumn{1}{c}{} & \multicolumn{1}{c}{\textbf{А}} & \multicolumn{1}{c}{\textbf{Б}} & \multicolumn{1}{c}{\textbf{В}} & \multicolumn{1}{c}{\textbf{Г}} & \multicolumn{1}{c}{\textbf{Д}} \\ \cline{2-6}
         \textbf{1} & \phantom{X} & \phantom{X} & \phantom{X} & \phantom{X} & \phantom{X} \\ \cline{2-6}
         \textbf{2} & & & & & \\ \cline{2-6}
         \textbf{3} & & & & & \\ \cline{2-6}
    \end{tabular}
    \endgroup
}
% --- ДОДАТКОВІ МАКРОСИ ЗБІРНИКА ---
\newcommand{\taskBlock}[1]{%
\par\vspace{0.45cm}
\noindent{\color{mainGreen}\rule{\linewidth}{0.8pt}}\par\vspace{0.12cm}
\noindent{\small\bfseries\color{mainGreen}#1}\par\vspace{0.25cm}
}
\newcounter{zad}
\newcommand{\zadtask}[1]{\stepcounter{zad}\noindent\textbf{\thezad.}\ \ #1\par\nopagebreak\vspace{0.2cm}}
\newcommand{\zadnum}{\stepcounter{zad}\textbf{\thezad.}\ \ }
\newcounter{ans}
\newcommand{\ansitem}[1]{\stepcounter{ans}\noindent\textbf{\theans.}\ #1\par\vspace{0.07cm}}
% мітка року: нерозривна, притиснута праворуч; якщо не вміщається -- праворуч на новому рядку
\newcommand{\nmtyear}[1]{\unskip\nobreak\finalhyphendemerits=0 \hfill\penalty100\hskip1em\hbox{}\nobreak\hfill{\small\color{yearOrange}\mbox{\rule{0pt}{2.3ex}(НМТ~#1)}}}
% --- МАКРОС КОРОТКОГО РОЗВ'ЯЗКУ ---
\newcommand{\solution}[1]{%
\par\vspace{0.12cm}
\begin{tcolorbox}[colback=titleBg!45, colframe=mainGreen!55, boxrule=0.4pt, arc=2pt,
    left=6pt, right=6pt, top=3pt, bottom=3pt, breakable]
\footnotesize\textbf{Короткий розв'язок.}\ #1
\end{tcolorbox}
\par\vspace{0.12cm}
}
% Розв'язки й відповіді (є до завдань НМТ-2026): \showsolutionstrue -- показати, \showsolutionsfalse -- сховати
\newif\ifshowsolutions
\showsolutionsfalse
% --- ЗАГОЛОВКИ З ПУНКТАМИ КЛІКАБЕЛЬНОГО ЗМІСТУ ---
\setcounter{tocdepth}{2}
\newcommand{\chapterTitle}[1]{%
\clearpage\phantomsection\addcontentsline{toc}{section}{#1}%
\sectionTitle{#1}}
\newcommand{\typeTitle}[1]{%
\needspace{6\baselineskip}%
\phantomsection\addcontentsline{toc}{subsection}{\quad #1}%
\par\vspace{0.3cm}
\noindent{\large\bfseries\color{yearOrange}#1}\par\vspace{0.08cm}
\noindent{\color{yearOrange}\rule{\linewidth}{0.4pt}}\par\nopagebreak\vspace{0.2cm}}
\raggedbottom
% усередині одного завдання сторінка не розривається: нерозривний вертикальний проміжок
% і нерозривна межа абзаців (умова ніколи не відділяється від варіантів, рисунка чи сітки)
\newcommand{\nbvspace}[1]{\nopagebreak\vspace{#1}}
\newcommand{\nmtnobreak}{\ifvmode\penalty10000 \fi}
% --- ЄДИНИЙ МАКЕТ ЗАВДАНЬ НА ВІДПОВІДНІСТЬ ---
% мітка в колонці сталої ширини, текст -- у parbox: продовження довгого рядка
% вирівнюється під текстом, а не під міткою; відступ між пунктами однаковий скрізь
\newlength{\nmtMatchLab}\setlength{\nmtMatchLab}{1.6em}
\newlength{\nmtMatchSep}\setlength{\nmtMatchSep}{0.28cm}
\newcommand{\matchHead}[1]{\par\noindent\textit{#1}\par\nopagebreak\vspace{0.2cm}}
% шаблонний заголовок стовпця зі старого макета -- лише коли завдання не має власного \matchHead
\makeatletter\newcommand{\nmtHeadUnless}[2]{\in@{\matchHead}{#1}\ifin@\else#2\fi}\makeatother
\newcommand{\matchItem}[2]{\par\noindent\makebox[\nmtMatchLab][l]{\textbf{#1}}%
\parbox[t]{\dimexpr\linewidth-\nmtMatchLab\relax}{\raggedright #2}\par\nopagebreak\vspace{\nmtMatchSep}}
\newcommand{\ansTheme}[1]{\par\vspace{0.25cm}\noindent{\bfseries\color{mainGreen}#1}\par\vspace{0.12cm}}
\newcommand{\ansType}[1]{\par\vspace{0.1cm}\noindent{\itshape\color{yearOrange}#1}\par\vspace{0.08cm}}

\graphicspath{{../1. Числа. Звичайні та десяткові дроби. Модуль числа/}{../10. Основи планіметрії/}{../11. Трикутники/}{../12. Паралелограм/}{../13. Прямокутник/}{../14. Квадрат/}{../15. Ромб/}{../16. Трапеція/}{../17. Коло, круг і їх елементи/}{../18. Координати і вектори на площині і в просторі. Рівняння кола/}{../19. Лінійні та квадратні нерівності. Системи лінійних нерівностей. Метод інтервалів/}{../2.відсотки, арифметичні задачі, подільність/}{../20. Функція та її властивості/}{../21.  Графіки елементарних функцій, їх побудова графіків функцій та їх перетворення/}{../22. Модульні  рівняння та нерівності/}{../23. Ірраціональні рівняння та нерівності/}{../24. Арифметична прогресія/}{../25. Геометрична прогресія/}{../26. Тригонометричні вирази/}{../27. Тригонометричні рівняння/}{../28. Показникова функція. Показникові рівняння/Показникові рівняння/}{../28. Показникова функція. Показникові рівняння/Показникові функція і вирази/}{../29. Показникові нерівності/}{../3. Степені. Одночлени. стандартний вигляд числа/}{../30. Логарифм. Логарифмічна функція/}{../31. Логарифмічні рівняння/}{../32. Логарифмічні нерівності/}{../33. Похідна/}{../34. Первісна/}{../35. Комбінаторика/}{../36. Теорія ймовірності/}{../37. Статистика/}{../38. Призма. Паралелепіпед. Куб/}{../39. Піраміда/}{../4.Многочлени. ФСМ/}{../40. Циліндр/}{../41. Конус/}{../42. Куля. Сфера/}{../43. Параметри/}{../5. дробові раціональні вирази і перетворення/}{../6. Лінійні рівняння/}{../7. Квадратні корені. Корені вищих степенів. Раціональний степінь/}{../8. квадратні рівняння, і рівняння, що до них зводяться/}{../9. Системи рівнянь/}}
\renewcommand{\instructionBox}[1]{%
\par\vspace{0.3cm}\noindent
\begin{tcolorbox}[nobeforeafter, width=\linewidth, colback=instrBg, colframe=instrBorder, boxrule=0.6pt, arc=2pt,
    left=8pt, right=8pt, top=4pt, bottom=4pt]
\centering\small\bfseries #1
\end{tcolorbox}
\par\nopagebreak\vspace{0.15cm}\nopagebreak
}
\renewcommand{\nmtAnswerBox}{%
\par\nopagebreak\vspace{0.2cm}\noindent
Відповідь:\ \
\begin{tabular}{@{}*{5}{|>{\centering\arraybackslash}p{0.8cm}}|@{\hspace{0.1cm}\textbf{,}\hspace{0.1cm}}*{3}{|>{\centering\arraybackslash}p{0.8cm}}|@{}}
\hline
\rule[-0.2cm]{0pt}{0.7cm} & & & & & & & \\
\hline
\end{tabular}
\par\vspace{0.3cm}
}

\begin{document}
% макроси бази, потрібні аналогам
% --- сумісність зі старою базою (діє, якщо тема не має власного визначення) ---
\providecommand{\trigStyle}[1]{#1}
\providecommand{\answerTableSmall}[5]{%
\begingroup\setlength{\tabcolsep}{2pt}\small\nmtSetCell
\begin{tabular}{|*{5}{>{\centering\arraybackslash}m{\nmtcellw}|}}
\hline
\rule[-0.2cm]{0pt}{0.6cm}\textbf{А} & \textbf{Б} & \textbf{В} & \textbf{Г} & \textbf{Д} \\
\hline
\rule[-0.4cm]{0pt}{0.9cm}#1 & \rule[-0.4cm]{0pt}{0.9cm}#2 & \rule[-0.4cm]{0pt}{0.9cm}#3 & \rule[-0.4cm]{0pt}{0.9cm}#4 & \rule[-0.4cm]{0pt}{0.9cm}#5 \\
\hline
\end{tabular}\endgroup}
\providecommand{\matchingLayout}[3]{%
    \noindent
    \begin{minipage}[t]{0.36\textwidth}\vspace{0pt}\raggedright
        #1
    \end{minipage}%
    \hfill
    \begin{minipage}[t]{0.43\textwidth}\vspace{0pt}\raggedright
        #2
    \end{minipage}%
    \hfill
    \begin{minipage}[t]{0.19\textwidth}
        \vspace{0pt}
        \begin{flushright}
        #3
        \end{flushright}
    \end{minipage}}
\providecommand{\matchingLayoutBottom}[2]{%
    \noindent
    \begin{minipage}[t]{0.48\textwidth}\vspace{0pt}\raggedright
        #1
    \end{minipage}%
    \hfill
    \begin{minipage}[t]{0.48\textwidth}\vspace{0pt}\raggedright
        #2
    \end{minipage}
    \vspace{0.5cm}
    \noindent
    \matchingGrid}
\providecommand{\answerListVertical}[5]{%
    \par\nopagebreak\vspace{0.2cm}
    \begin{itemize}[itemsep=0.4cm, leftmargin=1.5cm, labelsep=0.5cm]
        \item[\textbf{А}] #1
        \item[\textbf{Б}] #2
        \item[\textbf{В}] #3
        \item[\textbf{Г}] #4
        \item[\textbf{Д}] #5
    \end{itemize}
    \vspace{0.2cm}}
\setcounter{zad}{0}

\chapterTitle{Тренувальний варіант НМТ з математики № 4}

\noindent{\small Варіант складено за структурою НМТ 2023--2026: 15 завдань з вибором однієї відповіді, 3 на встановлення відповідності й 4 з короткою відповіддю (максимум 32 бали). Усі 22 завдання -- авторські, складені за зразком справжніх завдань НМТ 2023--2026 (такі самі за змістом і дистракторами). Відповіді -- на останній сторінці.}\par\vspace{0.4cm}

\instructionBox{Завдання 1--15 мають по п'ять варіантів відповіді, з яких лише \textbf{один правильний}. Виберіть правильний, на Вашу думку, варіант відповіді й позначте його в бланку відповідей.}

\begin{samepage}
\zadtask{Велосипедист проїхав $0{,}4$ запланованого маршруту. Після цього йому залишилося проїхати $36$~км. Визначте довжину всього маршруту.}
\answerTable{$14{,}4$~км}{$21{,}6$~км}{$50{,}4$~км}{$60$~км}{$90$~км}
\par\penalty-20
\end{samepage}

\begin{samepage}
\noindent
\begin{minipage}[c]{0.46\textwidth}
\zadnum На діаграмі відображено частку учнів (у~\%), які з першої спроби склали практичний іспит з водіння, у п'яти автошколах. Визначте автошколу, у якій ця частка становила $\dfrac{3}{5}$ від загальної кількості учнів цієї автошколи.
\end{minipage}
\hfill
\begin{minipage}[c]{0.52\textwidth}
\centering
\begin{nmtfit}\begin{tikzpicture}[x=0.9cm, y=0.045cm, thick]
    \foreach \y in {5,10,...,100} {
        \draw[gray!25, thin] (0,\y) -- (5.3,\y);
    }
    \foreach \y in {10,20,...,100} {
        \node[left, font=\scriptsize] at (0,\y) {\y};
    }
    \node[left, font=\scriptsize] at (0,0) {0};
    \draw[->] (0,0) -- (0,108);
    \draw[->] (0,0) -- (5.6,0);
    \node[rotate=90, font=\scriptsize] at (-0.95,50) {Частка учнів, \%};
    \filldraw[fill=mainGreen!75, draw=none] (0.55,0) rectangle ++(0.5,60);
    \filldraw[fill=yearOrange!75, draw=none] (1.55,0) rectangle ++(0.5,35);
    \filldraw[fill=gray!55, draw=none] (2.55,0) rectangle ++(0.5,75);
    \filldraw[fill=mainGreen!45, draw=none] (3.55,0) rectangle ++(0.5,30);
    \filldraw[fill=yearOrange!45, draw=none] (4.55,0) rectangle ++(0.5,40);
    \node[rotate=45, anchor=north east, font=\scriptsize, inner sep=2pt] at (0.95,-2) {«Маршрут»};
    \node[rotate=45, anchor=north east, font=\scriptsize, inner sep=2pt] at (1.95,-2) {«Старт»};
    \node[rotate=45, anchor=north east, font=\scriptsize, inner sep=2pt] at (2.95,-2) {«Кермо»};
    \node[rotate=45, anchor=north east, font=\scriptsize, inner sep=2pt] at (3.95,-2) {«Драйв»};
    \node[rotate=45, anchor=north east, font=\scriptsize, inner sep=2pt] at (4.95,-2) {«Світлофор»};
\end{tikzpicture}\end{nmtfit}
\end{minipage}
\par\nbvspace{0.5cm}
\answerTable{«Маршрут»}{«Старт»}{«Кермо»}{«Драйв»}{«Світлофор»}
\par\penalty-20
\end{samepage}

\begin{samepage}
\zadtask{На відрізку $MN$ завдовжки $34$~см позначено точку $P$ так, що відрізок $MP$ на $8$~см коротший за відрізок $PN$. Знайдіть довжину відрізка $PN$.}
\answerTable{$26$~см}{$13$~см}{$17$~см}{$25$~см}{$21$~см}
\par\penalty-20
\end{samepage}

\begin{samepage}
\zadtask{Розв'яжіть нерівність $\log_5(9 - 2x) < \log_5 4$.}
\answerTable{$(2{,}5;\, 4{,}5)$}{$(-\infty;\, 2{,}5)$}{$(2{,}5;\, +\infty)$}{$(4{,}5;\, +\infty)$}{$(-2{,}5;\, 4{,}5)$}
\par\penalty-20
\end{samepage}

\begin{samepage}
\zadtask{Скільки всього ребер має шестикутна піраміда?}
\answerTable{$6$}{$7$}{$12$}{$14$}{$18$}
\par\penalty-20
\end{samepage}

\begin{samepage}
\zadtask{Діаметр еритроцита (червоної клітини крові) людини становить приблизно $7{,}5$~\textit{мкм}, де $1$~\textit{мкм} (мікрометр)~--- одна мільйонна частина метра. Запишіть діаметр еритроцита (у~\textit{м}) у стандартному вигляді.}
\answerTable{$75 \cdot 10^{-7}$}{$7{,}5 \cdot 10^{6}$}{$0{,}75 \cdot 10^{-5}$}{$7{,}5 \cdot 10^{-6}$}{$7{,}5 \cdot 10^{-5}$}
\par\penalty-20
\end{samepage}

\begin{samepage}
\zadtask{Якому проміжку належить корінь рівняння $3^{x+2} - 3^{x} = \dfrac{8}{27}$?}
\answerTable{$(-\infty;\, -3)$}{$[-3;\, -2)$}{$[-2;\, 0)$}{$[0;\, 2)$}{$[2;\, +\infty)$}
\par\penalty-20
\end{samepage}

\begin{samepage}
\noindent
\begin{minipage}[c]{0.58\textwidth}
\zadnum На рисунку зображено графік функції $y = f(x)$, визначеної на проміжку $[-4;\, 4]$. Укажіть точку максимуму цієї функції.
\end{minipage}
\hfill
\begin{minipage}[c]{0.40\textwidth}
\centering
\begin{nmtfit}\begin{tikzpicture}[scale=0.5, thick]
    \draw[gray!40, thin] (-4.5,-2.5) grid (4.5,5.5);
    \draw[->, thick] (-4.8,0) -- (4.9,0) node[below] {$x$};
    \draw[->, thick] (0,-2.8) -- (0,5.9) node[left] {$y$};
    \node[below left] at (0,0) {$0$};
    \node[below] at (1,0) {$1$};
    \node[below left] at (-4,0) {$-4$};
    \node[below] at (4,0) {$4$};
    \node[right] at (0,1) {$1$};
    \draw[mainGreen!80!black, very thick] (-4,-1) .. controls (-3.3,1.5) and (-2.7,3) .. (-2,3)
        .. controls (-0.7,3) and (0.7,-2) .. (2,-2)
        .. controls (2.8,-2) and (3.6,2) .. (4,5);
    \fill[mainGreen!80!black] (-4,-1) circle (3pt);
    \fill[mainGreen!80!black] (4,5) circle (3pt);
    \node[fill=white, inner sep=1pt, font=\small] at (2.3,4.2) {$y = f(x)$};
\end{tikzpicture}\end{nmtfit}
\end{minipage}
\par\nbvspace{0.2cm}
\answerTable{$-2$}{$2$}{$3$}{$4$}{$5$}
\par\penalty-20
\end{samepage}

\begin{samepage}
\zadtask{У рівнобедреному трикутнику $ABC$ ($AB = BC$) $BH$~--- висота, проведена до сторони $AC$. Які з наведених тверджень є \textit{правильними}?\par
\nopagebreak\nbvspace{0.1cm}
\begin{enumerate}[label=\textbf{\Roman*.}, align=left, leftmargin=2.6em, labelsep=0.5em, itemsep=2pt, topsep=2pt]
\item $AH = HC$.
\item $BH > AB$.
\item $\angle ABH = \angle CBH$.
\end{enumerate}
\nopagebreak\nbvspace{0.1cm}}
\answerTable{лише I}{лише I та II}{лише I та III}{лише II та III}{I, II та III}
\par\penalty-20
\end{samepage}

\begin{samepage}
\zadtask{Укажіть інтеграл, який має \textit{найбільше} значення.}
\answerTableTall{$\displaystyle\int_0^1 8\,dx$}{$\displaystyle\int_1^4 3\,dx$}{$\displaystyle\int_5^6 4\,dx$}{$\displaystyle\int_0^6 1\,dx$}{$\displaystyle\int_{-2}^{1} 4\,dx$}
\par\penalty-20
\end{samepage}

\begin{samepage}
\zadtask{У вазі лежать $40$ цукерок, частина з яких~--- із шоколадною начинкою. Імовірність того, що навмання взята з вази цукерка виявиться із шоколадною начинкою, дорівнює $\dfrac38$. Скільки цукерок із шоколадною начинкою лежить у вазі?}
\answerTable{$5$}{$8$}{$15$}{$25$}{$120$}
\par\penalty-20
\end{samepage}

\begin{samepage}
\zadtask{У прямокутній системі координат у просторі задано конус із вершиною $S(3;\, -2;\, 6)$ і центром основи $Q(2;\, 2;\, -2)$. Визначте висоту конуса.}
\answerTable{$\sqrt{13}$}{$\sqrt{33}$}{$\sqrt{65}$}{$9$}{$81$}
\par\penalty-20
\end{samepage}

\begin{samepage}
\zadtask{Знайдіть значення виразу $3\sin\left(-\dfrac{3\pi}{2}\right) - \cos 5\pi$.}
\answerTable{$-4$}{$-2$}{$2$}{$3$}{$4$}
\par\penalty-20
\end{samepage}

\begin{samepage}
\noindent
\begin{minipage}[c]{0.58\textwidth}
\zadnum У прямокутній трапеції $ABCD$ ($AD \parallel BC$, $\angle A = \angle B = 90^\circ$) на бічній стороні $AB$ вибрано точку $K$ так, що $\angle CKD = 90^\circ$ (див. рисунок). $CK = 4\sqrt{3}$~см, $KD = 8$~см, $\angle BCK = 60^\circ$. Визначте площу трапеції $ABCD$.
\end{minipage}
\hfill
\begin{minipage}[c]{0.40\textwidth}
\centering
\begin{nmtfit}\begin{tikzpicture}[scale=0.36, thick, line join=round]
    \coordinate (A) at (0,0);
    \coordinate (B) at (0,10);
    \coordinate (C) at (3.464,10);
    \coordinate (D) at (6.928,0);
    \coordinate (K) at (0,4);
    \draw (A) -- (B) -- (C) -- (D) -- cycle;
    \draw (C) -- (K) -- (D);
    \pic[draw, angle radius=0.2cm] {right angle = D--A--B};
    \pic[draw, angle radius=0.2cm] {right angle = A--B--C};
    \pic[draw, angle radius=0.2cm] {right angle = D--K--C};
    \pic[draw, angle radius=0.45cm, pic text={\scriptsize $60^\circ$}, angle eccentricity=1.8] {angle = B--C--K};
    \fill (K) circle (3pt);
    \node[below left] at (A) {$A$};
    \node[above left] at (B) {$B$};
    \node[above right] at (C) {$C$};
    \node[below right] at (D) {$D$};
    \node[left] at (K) {$K$};
\end{tikzpicture}\end{nmtfit}
\end{minipage}
\par\nbvspace{0.2cm}
\answerTable{$16\sqrt{3}$~см$^2$}{$18\sqrt{3}$~см$^2$}{$30\sqrt{3}$~см$^2$}{$40\sqrt{3}$~см$^2$}{$60\sqrt{3}$~см$^2$}
\par\penalty-20
\end{samepage}

\begin{samepage}
\zadtask{Скільки всього коренів має рівняння $x^4 = 9x^2$?}
\answerTable{жодного}{один}{два}{три}{більше трьох}
\par\penalty-20
\end{samepage}

\instructionBox{У завданнях 16--18 до кожного з трьох рядків інформації, позначених цифрами, доберіть один правильний, на Вашу думку, варіант, позначений буквою.}

\begin{samepage}
\zadtask{Установіть відповідність між виразом \mbox{(1--3)} та його значенням \mbox{(А--Д)}.
\par\nopagebreak\nbvspace{0.3cm}
\noindent
\matchingLayout{
\matchHead{Вираз}
\matchItem{1}{$\left|\dfrac{3}{4} - 1{,}5\right|$}
\matchItem{2}{$\left(\dfrac{1}{2}\right)^{-3} \cdot 4^{-1}$}
\matchItem{3}{$\sqrt{3} \cdot \operatorname{tg}\dfrac{\pi}{6}$}
}{
\matchHead{Значення виразу}
\matchItem{А}{$-\dfrac{3}{4}$}
\matchItem{Б}{$\dfrac{3}{4}$}
\matchItem{В}{$1$}
\matchItem{Г}{$2$}
\matchItem{Д}{$3$}
}{
\matchingGrid
}}
\par\penalty-20
\end{samepage}

\begin{samepage}
\zadtask{До кожного початку речення \mbox{(1--3)} доберіть його закінчення \mbox{(А--Д)} так, щоб утворилося правильне твердження.
\par\nopagebreak\nbvspace{0.3cm}
\noindent
\matchingLayout{
\matchHead{Початок речення}
\matchItem{1}{Графіки функцій $y = \dfrac{4}{x}$ та $y = -x$}
\matchItem{2}{Графіки функцій $y = \sqrt{x}$ та $y = x - 2$}
\matchItem{3}{Графік функції $y = x^2 - 2$ та графік рівняння $x^2 + y^2 = 4$}
}{
\matchHead{Закінчення речення}
\matchItem{А}{не мають спільних точок.}
\matchItem{Б}{мають єдину спільну точку.}
\matchItem{В}{мають лише дві спільні точки.}
\matchItem{Г}{мають лише три спільні точки.}
\matchItem{Д}{мають чотири спільні точки.}
}{
\matchingGrid
}}
\par\penalty-20
\end{samepage}

\begin{samepage}
\zadtask{На рисунку зображено ромб $ABCD$, рівносторонній трикутник $CDK$ і прямокутний трикутник $AKN$ ($\angle AKN = 90^\circ$), що лежать в одній площині. Точки $A$, $D$ і $K$ лежать на одній прямій, катет $KN$ дорівнює меншій діагоналі ромба. Установіть відповідність між геометричною фігурою \mbox{(1--3)} та її площею \mbox{(А--Д)}, якщо $AK = 12$~см.
\par\nopagebreak\nbvspace{0.3cm}
\begin{center}
\begin{nmtfit}\begin{tikzpicture}[scale=0.4, thick, line join=round]
    \coordinate (A) at (0,0);
    \coordinate (D) at (6,0);
    \coordinate (K) at (12,0);
    \coordinate (B) at (3,5.196);
    \coordinate (C) at (9,5.196);
    \coordinate (N) at (12,-6);
    \filldraw[fill=yearOrange!20, draw=black] (A) -- (B) -- (C) -- (D) -- cycle;
    \filldraw[fill=mainGreen!20, draw=black] (D) -- (C) -- (K) -- cycle;
    \filldraw[fill=gray!20, draw=black] (A) -- (K) -- (N) -- cycle;
    \draw[dashed, thin] (B) -- (D);
    \draw[thin] (11.3,0) -- (11.3,-0.7) -- (12,-0.7);
    \node[left] at (A) {$A$};
    \node[above left] at (B) {$B$};
    \node[above right] at (C) {$C$};
    \node[below] at (D) {$D$};
    \node[right] at (K) {$K$};
    \node[below] at (N) {$N$};
\end{tikzpicture}\end{nmtfit}
\end{center}
\nopagebreak\nbvspace{0.2cm}
\noindent
\matchingLayout{
\matchHead{Геометрична фігура}
\matchItem{1}{трикутник $CDK$}
\matchItem{2}{трикутник $AKN$}
\matchItem{3}{ромб $ABCD$}
}{
\matchHead{Площа фігури}
\matchItem{А}{$36$~см$^2$}
\matchItem{Б}{$72$~см$^2$}
\matchItem{В}{$9\sqrt{3}$~см$^2$}
\matchItem{Г}{$18\sqrt{3}$~см$^2$}
\matchItem{Д}{$36\sqrt{3}$~см$^2$}
}{
\matchingGrid
}}
\par\penalty-20
\end{samepage}

\instructionBox{Розв'яжіть завдання 19--22. Одержані числові відповіді запишіть у бланку відповідей. Відповідь записуйте лише десятковим дробом, урахувавши положення коми. Знак «мінус» записуйте перед першою цифрою числа.}

\begin{samepage}
\zadtask{Знайдіть \textit{значення $x_0$}, при якому функція $y = -2x^4 + 25x^2 - 4$ набуває \textit{найбільшого} значення на проміжку $[-1;\, 4]$.}
\nmtAnswerBox
\par\penalty-20
\end{samepage}

\begin{samepage}
\zadtask{На шкільному концерті виступатимуть $7$ учасників: $3$ співаки та $4$ танцюристи (кожен учасник виступає один раз). Скільки всього існує різних варіантів порядку виступів учасників, якщо всі $3$ співаки мають виступати поспіль (один за одним)?}
\nmtAnswerBox
\par\penalty-20
\end{samepage}

\begin{samepage}
\noindent
\begin{minipage}[c]{0.58\textwidth}
\zadnum Радіус основи й висота конуса відносяться як $3 : 4$, а площа його бічної поверхні дорівнює $60\pi$. Площа повної поверхні півкулі дорівнює $48\pi$. У скільки разів об'єм конуса більший за об'єм півкулі?
\end{minipage}
\hfill
\begin{minipage}[c]{0.40\textwidth}
\centering
\begin{nmtfit}\begin{tikzpicture}[scale=0.85, thick, line join=round]
    \draw[dashed] (-1.2,0) arc (180:0:1.2 and 0.34);
    \draw (-1.2,0) arc (180:360:1.2 and 0.34);
    \draw (-1.2,0) -- (0,1.6) -- (1.2,0);
    \begin{scope}[shift={(3.1,0)}]
        \draw (-0.8,0) arc (180:0:0.8);
        \draw[dashed] (-0.8,0) arc (180:0:0.8 and 0.23);
        \draw (-0.8,0) arc (180:360:0.8 and 0.23);
    \end{scope}
\end{tikzpicture}\end{nmtfit}
\end{minipage}
\par\nbvspace{0.2cm}
\nmtAnswerBox
\par\penalty-20
\end{samepage}

\begin{samepage}
\zadtask{Знайдіть суму всіх цілих значень $a$, за кожного з яких корінь рівняння
\[
\dfrac{\log_2(2x + a + 1) - 4}{\sqrt{x^2 - a^2}} = 0
\]
є цілим числом.}
\nmtAnswerBox
\par\penalty-20
\end{samepage}

\clearpage
\sectionTitle{Відповіді}

\begin{center}\small\renewcommand{\arraystretch}{1.25}
\begin{tabular}{|c|c|p{0.62\textwidth}|}\hline
\textbf{№} & \textbf{Відповідь} & \textbf{Джерело завдання} \\ \hline
1 & Г & аналог (авторське завдання за зразком НМТ) \\ \hline
2 & А & аналог (авторське завдання за зразком НМТ) \\ \hline
3 & Д & аналог (авторське завдання за зразком НМТ) \\ \hline
4 & А & аналог (авторське завдання за зразком НМТ) \\ \hline
5 & В & аналог (авторське завдання за зразком НМТ) \\ \hline
6 & Г & аналог (авторське завдання за зразком НМТ) \\ \hline
7 & Б & аналог (авторське завдання за зразком НМТ) \\ \hline
8 & А & аналог (авторське завдання за зразком НМТ) \\ \hline
9 & В & аналог (авторське завдання за зразком НМТ) \\ \hline
10 & Д & аналог (авторське завдання за зразком НМТ) \\ \hline
11 & В & аналог (авторське завдання за зразком НМТ) \\ \hline
12 & Г & аналог (авторське завдання за зразком НМТ) \\ \hline
13 & Д & аналог (авторське завдання за зразком НМТ) \\ \hline
14 & В & аналог (авторське завдання за зразком НМТ) \\ \hline
15 & Г & аналог (авторське завдання за зразком НМТ) \\ \hline
16 & 1~--~Б, 2~--~Г, 3~--~В & аналог (авторське завдання за зразком НМТ) \\ \hline
17 & 1~--~А, 2~--~Б, 3~--~Г & аналог (авторське завдання за зразком НМТ) \\ \hline
18 & 1~--~В, 2~--~А, 3~--~Г & аналог (авторське завдання за зразком НМТ) \\ \hline
19 & $2{,}5$ & аналог (авторське завдання за зразком НМТ) \\ \hline
20 & $720$ & аналог (авторське завдання за зразком НМТ) \\ \hline
21 & $2{,}25$ & аналог (авторське завдання за зразком НМТ) \\ \hline
22 & $-45$ & аналог (авторське завдання за зразком НМТ) \\ \hline
\end{tabular}\end{center}

\noindent{\small\color{gray!80!black}Оцінювання: завдання 1--15 -- по 1 балу; 16--18 -- по 1 балу за кожну правильно встановлену відповідність (до 3 балів); 19--22 -- по 2 бали. Разом 32 бали.}

\end{document}
